\documentclass[11pt]{article}
\usepackage[T1]{fontenc}
\usepackage[utf8]{inputenc}
\usepackage[a4paper,margin=2.2cm]{geometry}
\usepackage{amsmath,amssymb}
\pagestyle{plain}
\setlength{\parskip}{4pt}
\newcommand{\A}{\mathcal A}
\newcommand{\G}{\mathcal G}

\begin{document}

\begin{center}
{\Large\bfseries Object Provenance for the Orbit-Space Gap Chain}\\[4pt]
{\normalsize A common page, reproduced verbatim in
ai.viXra:2602.0033, 2602.0035 and 2602.0036}\\[4pt]
Lluis Eriksson \quad\textbullet\quad July 2026
\end{center}

\noindent\rule{\textwidth}{0.8pt}

\vspace{2pt}
\noindent
\textbf{Why this page exists.} Four papers of this series state, prove, or consume
spectral claims about objects that are routinely called ``the gap on the orbit
space''. They are not the same object. A pointwise curvature \emph{claim} stated in
one of them was consumed, in three places, as if it were established and as if it
applied to the same object. It is neither: its printed proof is not presently
established, and the sound synthetic result available elsewhere concerns a measure and
an operator that \emph{have not been identified} with the ones the consuming steps
need. The error was invisible because every paper used
the same word. This page fixes the dictionary. \textbf{A citation discharges a
hypothesis only when all four columns agree.}

\section*{Notation barrier}

\textbf{This page writes $\nu$ for $\nu_{\rm push}:=\pi_\#\mathrm{Vol}_{\A}$, and never
for anything else.} The host papers of this chain use the same letter for a different
object: \texttt{2602.0033} (Theorem 3.1, Step 3) and \texttt{2602.0035} (Theorem 5.3,
Step 2) both write ``the uniform measure $\nu$'' for the normalised Riemannian volume
\[
\nu_{\rm Vol}:=d\mathrm{Vol}_B/\mathrm{Vol}_B(B).
\]
\textbf{$\nu_{\rm push}$ and $\nu_{\rm Vol}$ are not identified}, and identifying them
is exactly what (I2b) below forbids without a statement about the orbit-volume density.
A reader who carries the host's $\nu$ into this page, or this page's $\nu$ into the
host's Step 3, will read a printed derivation as though it already used the corrected
measure. \emph{Where the distinction bites, the subscripts are written out.}

\section*{The dictionary}

\begin{center}
\footnotesize
\setlength{\tabcolsep}{4pt}
\begin{tabular}{|l|l|l|l|}
\hline
\textbf{Source} & \textbf{Space, measure} & \textbf{Operator} & \textbf{What it can conclude}\\
\hline
\texttt{2602.0036}
 & $B_{\rm reg}$, $d\mathrm{Vol}_B$
 & $\mathcal K_B^{\rm reg}$: unweighted
 & \emph{targets} a pointwise Ricci\\
 & & Laplace--Beltrami \emph{expression};
 & bound on the regular stratum;\\
 & & global realisation unspecified
 & printed proof not presently\\
 & & & established\\
\hline
\texttt{2602.0046}
 & $B$, $\nu=\pi_\#\mathrm{Vol}_{\A}$
 & $K_\nu\ge0$, operator of the
 & global Poincar\'e/LSI,\\
 & & form on $L^2(B,\nu)$ & singular locus included\\
\hline
\texttt{2602.0033},
 & spatial orbit space; printed
 & weighted diffusion of
 & neither a Poincar\'e/LSI\\
as printed
 & \emph{regular-stratum formula}
 & $\mu_{\rm print}$, \emph{if} a global
 & statement nor a Hamiltonian\\
 & $d\mu_{\rm print}=Z_{\rm print}^{-1}e^{-S_{\rm eff}}d\mathrm{Vol}_B$
 & closed form is constructed
 & gap is presently\\
 & on $B_{\rm reg}$, \textbf{against $\mathrm{Vol}_B$, not $\nu$};
 &
 & established\\
 & \emph{no global extension specified} & & \\
\hline
route (A):
 & $(B,\nu)$; a candidate \emph{global}
 & weighted form
 & needs $w$ globally, and\\
keep $\mu_{\rm print}$
 & \emph{extension} of $\mu_{\rm print}$,
 & over $K_\nu$
 & $\operatorname{osc}(S_{\rm eff}+\log w)$\\
 & exponent $S_{\rm eff}+\log w$ & & \\
\hline
route (B):
 & $(B,\nu)$; a \emph{separately}
 & weighted form
 & needs a theorem that\\
new measure
 & \emph{defined candidate} measure
 & over $K_\nu$
 & this is the right object\\
 & $d\widetilde\mu=\widetilde Z^{-1}e^{-S_{\rm eff}}d\nu$ & & \\
\hline
\texttt{2602.0035},
 & $L^2(B,d\mathrm{Vol}_B)$, global
 & Hamiltonian on
 & parts (i) and (ii) are not\\
as printed
 & realisation not fixed;
 & $L^2(B,d\mathrm{Vol}_B)$;
 & established by the printed\\
 & $\mu_{0,\rm print}
 =\psi_{0,\rm print}^2\,d\mathrm{Vol}_B$
 & diffusion of $\mu_{0,\rm print}$
 & proofs\\
\hline
candidate
 & $L^2(\A)^{\G}\cong L^2(B,\nu)$;
 & $H_\nu=\frac{g^2}{2}K_\nu
 +V_{\rm pot}$;
 & qualitative fixed-volume\\
physical
 & $\mu_{0,\nu}=\psi_{0,\nu}^2\,\nu$
 & then diffusion of
 & positivity; a rate only\\
restatement
 &
 & $\mu_{0,\nu}$
 & under $(\text{A-osc})_\nu$\\
\hline
\end{tabular}
\end{center}

\section*{The four identifications that must never be silent}

\paragraph{(I1) $\;L^2(B,\nu)\cong L^2(\A,\mathrm{Vol}_{\A})^{\G}$ --- exact, and free.}
Pullback $f\mapsto f\circ\pi$ gives
$\int_B|f|^2\,d\nu=\int_{\A}|f\circ\pi|^2\,d\mathrm{Vol}_{\A}$, which is the
definition of $\nu=\pi_\#\mathrm{Vol}_{\A}$. This is an isometry, not an
approximation, and it holds across the singular locus. It is the reason the
metric-measure quotient is the natural physical object: gauge-invariant functions on
the smooth compact $\A$ \emph{are} $L^2(B,\nu)$.

\paragraph{(I2) $\;K_\nu$ versus $\mathcal K_B^{\rm reg}$ --- not free, and not even
the same kind of object.}
Fix the convention here, so that nothing depends on the host paper's. On the regular
stratum write $d\nu=w\,d\mathrm{Vol}_B$, with $w$ the orbit-volume density, and define
the \emph{formal differential expressions}
\[
\mathcal L_B^{\rm reg}:=\operatorname{div}_B\nabla,
\qquad
\mathcal L_\nu^{\rm reg}:=w^{-1}\operatorname{div}_B(w\nabla),
\qquad
\mathcal K_\bullet^{\rm reg}:=-\mathcal L_\bullet^{\rm reg},
\]
so that
\[
\mathcal L_\nu^{\rm reg} f=\mathcal L_B^{\rm reg} f+\langle\nabla\log w,\nabla f\rangle .
\]
Let $K_\nu\ge0$ denote the \emph{global} self-adjoint operator associated with the
closed Dirichlet form on $L^2(B,\nu)$.

\textbf{No global self-adjoint realisation of $\mathcal K_B^{\rm reg}$ is fixed by
this page.} $B_{\rm reg}$ is the regular stratum with the singular locus removed and
is in general an incomplete manifold; $\mathcal K_B^{\rm reg}$ on
$C_c^\infty(B_{\rm reg})$ is symmetric and non-negative, but it has not been shown to
be essentially self-adjoint, and it may admit several self-adjoint extensions.
Choosing the Friedrichs extension is already a decision about behaviour across
$B\setminus B_{\rm reg}$, and must be specified and justified separately --- this is
one of the gaps that forces the withdrawal of the global corollary in
\texttt{2602.0036}.

The two \emph{regular-stratum differential expressions} coincide only if $w$ is
constant on each relevant component. \emph{Constant stabiliser type does not
by itself imply constant orbit volume} --- on the principal stratum all stabilisers
are already conjugate, and the induced metric on the orbits, hence their volume, may
still vary with the configuration; reducible connections make a global identification
less tenable still. \textbf{A spectral bound for the global $K_\nu$ does not by itself
transfer to any chosen self-adjoint realisation of $\mathcal K_B^{\rm reg}$, or
conversely.} Comparison theorems under further hypotheses are not excluded; what does
not exist is a free identification --- and on the $\mathcal K_B^{\rm reg}$ side, there
is not yet a single global operator for such a theorem to be about.

\paragraph{(I2b) $\;d\mathrm{Vol}_B$ versus $\nu$ \emph{at the level of the perturbed
measure} --- also not free.}
This is a separate step from (I2), and it is the one most easily skipped, because both
measures are written with the same exponent. On $B_{\rm reg}$,
$d\nu=w\,d\mathrm{Vol}_B$ with $w$ the orbit-volume density, so the measure a paper
prints against $\mathrm{Vol}_B$ is, written against $\nu$,
\[
d\mu_{\rm print}=Z_{\rm print}^{-1}e^{-S_{\rm eff}}\,d\mathrm{Vol}_B
=Z_{\rm print}^{-1}e^{-(S_{\rm eff}+\log w)}\,d\nu
\qquad\text{on }B_{\rm reg},
\]
an \emph{exact} rewriting of one and the same measure, with the same normaliser. The
candidate $\nu$-based measure $d\widetilde\mu=\widetilde Z^{-1}e^{-S_{\rm eff}}d\nu$
carries the exponent $S_{\rm eff}$ and, in general, a \emph{different} normaliser.
\textbf{On $B_{\rm reg}$ the printed measure has exponent $S_{\rm eff}+\log w$
relative to $\nu$, whereas the candidate $\nu$-based measure has exponent
$S_{\rm eff}$. They coincide only under an additional constancy or identification
statement for $w$, none of which is established here. Replacing $d\mathrm{Vol}_B$ by
$d\nu$ under a fixed exponent is therefore not a free rewriting.}

\emph{And the converse is not asserted either.} This page does not claim the two
measures are \emph{proved} distinct: that would need $w$ non-constant, up to
normalisation, on the relevant domain. Reducible connections make a gratuitous global
identification implausible, but implausible is not proved, and the discipline this
page enforces cuts both ways --- ``no identification has been proved'' must not become
``an inequality has been proved''.

\emph{Two routes are open, and they are not the same route.}
\textbf{(A) Keep the printed measure.} Put $\Phi_{\rm print}=S_{\rm eff}+\log w$, so
that $d\mu_{\rm print}=Z_{\rm print}^{-1}e^{-\Phi_{\rm print}}d\nu$. This needs a
global construction of the measure from the regular-stratum formula, control of $w$
approaching the singular locus, and a bound on
$\operatorname{osc}(S_{\rm eff}+\log w)$ --- for which a bound on
$\operatorname{osc}(S_{\rm eff})$ alone does not suffice: if $w$ degenerates at strata
of enlarged stabiliser, $\log w$ is exactly the term that destroys a
bounded-perturbation argument. \textbf{No such control is supplied anywhere in this
chain.}
\textbf{(B) Change to the $\nu$-natural measure} $\widetilde\mu$, which is
\emph{a separately defined candidate measure} and not the one the paper printed. Here
\texttt{2602.0046} can supply the base functional inequality for $K_\nu$, but a new
theorem is required to show that $\widetilde\mu$, and not $\mu_{\rm print}$, is the
correct Euclidean or transfer-theoretic object. Note that route (A) is not free
either: ``keep the printed measure'' still means \emph{constructing a global extension
of a regular-stratum formula}, which is itself one of the open obligations. \textbf{A $\nu$-based base inequality is therefore an input to one of these
two routes, not a repair of a printed $\mathrm{Vol}_B$-based derivation by change of
reference measure.} Neither route by itself supplies the four-dimensional-action to
three-dimensional-space step, the identification with $\mu_0$, or a transfer-operator
construction.

\paragraph{(I3) $\;\mu_{\rm print}$ has not been identified with $\mu_0=\psi_0^2\nu$
--- not free, and the load-bearing gap.}
With $K_\nu\ge0$ as in (I2) --- equivalently, by (I1), the operator of the kinetic
form $\mathcal E_\nu(f,f)=\int_B|\nabla f|^2\,d\nu$ on
$L^2(\A,\mathrm{Vol}_{\A})^{\G}$ --- \textbf{let $V_{\rm mult}$ be a specified real
multiplication potential} and set
$H_{\rm mult}=\frac{g^2}{2}K_\nu+V_{\rm mult}$, with a normalised, strictly positive
ground state $\psi_0$, and put $d\mu_0=\psi_0^2\,d\nu$. \emph{The subscripted name is
deliberate. This page is reproduced verbatim inside papers that have already bound
$V$ to a gauge-fixed potential and $W$ to the Weyl group, so no bare letter is safe
here: a page written to prevent symbol collisions must not cause one.} The
ground-state transformation \emph{then} gives
\[
E_1-E_0=\tfrac{g^2}{2}\,\lambda_1(\mu_0),\qquad f_0=-\log\psi_0^2 ,
\]
where $E_1$ denotes the first spectral value strictly above the simple ground-state
energy $E_0$, and, \emph{defined explicitly to avoid the reciprocal ambiguity in the
term ``Poincar\'e constant''},
\[
\lambda_1(\mu_0):=\inf\Bigl\{\mathcal E_{\mu_0}(f,f)\;:\;
f\in\operatorname{Dom}\mathcal E_{\mu_0},\ \textstyle\int f\,d\mu_0=0,\
\|f\|_{L^2(\mu_0)}=1\Bigr\}.
\]
That is, $\lambda_1(\mu_0)$ is the Poincar\'e \emph{spectral gap} --- the largest
$\rho\ge0$ with $\operatorname{Var}_{\mu_0}(f)\le\rho^{-1}\mathcal E_{\mu_0}(f,f)$ for
all admissible $f$ --- and \emph{not} the constant $C_P=\rho^{-1}$ of the variance
inequality, which some authors also call the Poincar\'e constant.

The Gibbs measure of a Euclidean effective action is \emph{a priori} a distinct
mathematical object, and its weighted Dirichlet-form gap is \emph{a priori} a distinct
spectral quantity. In the symbols fixed above,
\[
\boxed{\;\text{No globally defined printed or corrected Euclidean measure has been
identified with }\mu_0=\psi_0^2\,\nu.\;}
\]
That covers $\mu_{\rm print}$, route (A)'s global extension of it, and route (B)'s
$\widetilde\mu$ alike --- stated this way so that (I3) neither singles out one
candidate nor silently presupposes (I2b). No such
identification has been proved here; in special models the objects may coincide, but
passing from one to the other in general requires a transfer-operator or
Osterwalder--Seiler identification theorem, not a change of notation. \textbf{An
oscillation bound for $S_{\rm eff}$ does not by itself bound
$\operatorname{osc}(-\log\psi_0^2)$.}

\section*{Constants travel with their metric normalisation}

$\operatorname{Ric}\ge N_c/2$ for $\langle X,Y\rangle=-\operatorname{tr}(XY)$ and
$\operatorname{Ric}\ge N_c/4$ for $-2\operatorname{tr}(XY)$ are the \emph{same} bound
in two normalisations, not a sharpening of one by the other. Any statement quoting
$N_c/2$ or $N_c/4$ must name its inner product.

\section*{How to use this page}

Before writing that a companion result discharges a hypothesis, fill the four columns
for the result and for the hypothesis. If any column differs, the citation does not
discharge; at most it supplies an input to a further identification, which must then
be stated as its own step. In particular:

\begin{itemize}
\item \texttt{2602.0046} supplies global geometry on the metric-measure object
identified by (I1): a bound for $K_\nu$ on $L^2(B,\nu)$. It does not supply (I1) ---
that is free --- and it does \emph{not} supply a bound for any realisation of the
unweighted $\mathcal K_B^{\rm reg}$; nor does it attempt pointwise bounds on
$B_{\rm reg}$, as its own text says. Note that it reaches its conclusion by synthetic
quotient stability and \emph{explicitly bypasses the O'Neill computation}. That is now
the only sound geometric input in this chain: the O'Neill trace printed in
\texttt{2602.0036} omits the mixed horizontal--vertical sectional terms and does not
establish its own conclusion (see that record's erratum). \textbf{No pointwise Ricci
bound on $B_{\rm reg}$ may be cited as available.}
\item \texttt{2602.0033}, as printed, targets a weighted diffusion for
$d\mu_{\rm print}=Z_{\rm print}^{-1}e^{-S_{\rm eff}}d\mathrm{Vol}_B$ on
$B_{\rm reg}$ --- against $\mathrm{Vol}_B$, \emph{not} against $\nu$, and with no
global extension or associated closed Dirichlet form specified. It does not presently
establish a Euclidean Poincar\'e/LSI statement, and it does not establish a
Hamiltonian or transfer-operator gap without (I3). \textbf{Citing \texttt{2602.0046}
does not repair it by writing the same exponent against $\nu$}: by (I2b) that is not a
free rewriting. Either route (A) is taken and
$\operatorname{osc}(S_{\rm eff}+\log w)$ is controlled, or route (B) is taken and the
new measure is shown to be the right object.
\item \texttt{2602.0035} v2 \textbf{already uses a ground-state measure}, so
\textbf{(I3) is not the missing bridge in its printed quantitative proof} --- that was
the correction its v2 made to its own v1. Its printed-to-physical transition
\textbf{crosses (I2) and (I2b)}: the unweighted realisation against $K_\nu$, and
$d\mathrm{Vol}_B$ against $\nu$. The candidate physical restatement then requires
$(\text{A-osc})_\nu$, a hypothesis about \emph{its own} ground state.
\textbf{Identification (I3) re-enters only if that Hamiltonian construction is used to
identify a Euclidean or transfer-operator object}, as in the route towards (H1) of
\texttt{2602.0033}. Among the four identifications on this page, its qualitative
fixed-volume positivity needs only (I1); it additionally uses the standard
compact-elliptic and positivity-improving inputs --- self-adjoint realisation, compact
resolvent, connectedness --- and it does not require any treatment of the quotient
singularities. \emph{An earlier version of this page said flatly that the paper needs
(I3) for its quantitative bound. That was wrong, and it is the kind of error this page
exists to catch: it named the wrong hypothesis for a citation to discharge.}
\end{itemize}

\vspace{4pt}
\noindent\rule{\textwidth}{0.8pt}

\noindent\footnotesize
\textbf{Provenance.} This page was written after an audit found the same curvature
\emph{claim} consumed in three places under three different pairings of measure and
operator, and after a proposed repair --- citing the metric-measure result in place of
the regular-stratum one --- was found to change the object rather than fix the proof.
A later pass found that the regular-stratum claim is not established by its own
printed proof either, so the defect is now twofold: an invalid derivation in one
paper, and a correct theorem in another about a different object.
A third pass found the drift \emph{on this page}: the row for \texttt{2602.0033} had
recorded the printed measure against $\nu$ when the paper writes it against
$\mathrm{Vol}_B$, which is precisely the substitution (I2b) now forbids. It was
introduced by writing the row from the intended repair rather than from the printed
text. Hence the rule below, and hence (I2b) as a numbered step rather than a remark.
The same pass then caught the correction overshooting: (I2b) had asserted the two
measures \emph{are} different, which needs $w$ non-constant and is not proved here.
Both directions are now stated at the strength they have.
A fourth pass found that (I3) wrote its Hamiltonian as $\frac{g^2}{2}K_\nu+V$ --- and
$V$ is already bound, in \texttt{2602.0035}, to the gauge-fixed
$S_{\rm YM}-\log\det M$. A page reproduced verbatim inside that paper cannot borrow a
symbol the host has spent. The first repair reached for a bare $W$, which that same
paper binds to the Weyl group $S_{N_c}$: \textbf{no bare letter is safe in a page
meant to be embedded, only a subscripted name}. The same pass found the collision that
matters most, because it is in the \emph{measure} column: both host papers write
``the uniform measure $\nu$'' for the normalised $d\mathrm{Vol}_B$, which is what this
page's own (I2b) forbids identifying with $\pi_\#\mathrm{Vol}_{\A}$. Hence the notation
barrier above. A fifth pass found the row for
\texttt{2602.0035} built from its intended repair rather than its printed text --- the
identical error this page had already corrected once, for \texttt{2602.0033} --- and
found the claim that that paper needs (I3) to be simply wrong: its v2 already uses a
ground-state measure.
It records a failure mode that is not covered by the usual distinction between a false
statement and an invalid proof: \textbf{a correct theorem about the wrong object for
the consuming dependency}. Reproduce it verbatim; do not paraphrase it, because
paraphrase is how the columns drift.

\end{document}
